\begin{corollary}[Collapse rate as correctness approaches one]
\label{cor:collapse-error-rate}
Under Theorem~\ref{thm:shared-correctness}, fix finite $\lambda\ge0$
and suppose the initial correct-mode maximum is unique. For every other
correct mode $c$,
\begin{equation}
 \lim_{P\uparrow1}\frac{\log q_c^{(\lambda)}(P)}{\log(1-P)}
 =\frac{1}{\lambda+1+1/n}.
 \label{eq:collapse-error-exponent}
\end{equation}
Thus $q_c^{(\lambda)}(P)=(1-P)^{1/(\lambda+1+1/n)+o(1)}$:
larger fixed $\lambda$ reduces the decay exponent, but every minority mode
still vanishes. This is a logarithmic asymptotic relation, not an exact
finite-accuracy power law.
\end{corollary}

\begin{proof}
Theorem~\ref{thm:shared-correctness} gives $q\to\mathbf e_{c^\star}$,
so $S_q\to1$. The incorrect conditional flow in
Equation~\eqref{eq:shared-correctness-clock} preserves ordering and makes
$r_{\min}$ nondecreasing. For $h=\log(r_{\max}/r_{\min})$,
\[
 h'=-(r_{\max}-r_{\min})
 =-r_{\min}(e^h-1)\le-r_{\min}(0)h.
\]
Hence $r\to(1/n,\ldots,1/n)$ and $S_r\to1/n$, including trivially $n=1$.
Integrating $d\log q_c/d\tau=q_c-S_q$ and the correctness-clock equation
and taking time averages yields
\[
 \frac{\log q_c(\tau)}{\tau}\to-1,
 \qquad
 \frac{\log(P(\tau)/(1-P(\tau)))}{\tau}
 \to\lambda+1+1/n.
\]
Since $\log(1-P)=\log P-\log(P/(1-P))$ and $\log P\to0$, their ratio proves
Equation~\eqref{eq:collapse-error-exponent}.
\end{proof}

\paragraph{Correctness alone does not imply collapse.}
The independent-logit component matters. With purely shared geometry
$\mathsf K=\lambda vv^\top$, $\lambda>0$, all correct logits move equally,
so $q(t)=q(0)$ while
$\dot P=\lambda c_G(P)P^2(1-P)^2>0$ and $P\to1$.
For $\mathsf K=aI+bvv^\top$ with fixed $a>0$ and $b\ge0$, time rescaling
instead recovers Theorem~\ref{thm:shared-correctness} with $\lambda=b/a$.
Thus persistent mode-specific amplification, not increasing correctness by
itself, drives this collapse result.
